\documentclass[10pt,a4paper]{amsart}
\usepackage{amsmath,amssymb,amsthm,graphicx}
% Retained preamble for the Stacks Project chapter divisors.tex (Divisors).
% The chapter file pinned-original-source.tex carries only \input{preamble} on its first line, so no
% preamble text lives inside the chapter itself.  The Stacks Project uses one single shared preamble for
% every chapter; the snapshot stages it verbatim in sources/stacks-chow/preamble.tex.  The block below is
% copied mechanically, line for line, from that file's lines 14-18 (the Xy-pic setup) and lines 154-217
% (the native theorem declarations together with the mathematical macros).  Nothing was added, removed or
% redefined.  The wrapper class and the amsmath/amssymb/amsthm/graphicx base packages are supplied by the
% assembler, so the class/\externaldocument part of the shared preamble is not repeated here.
% For commutative diagrams we use Xy-pic
\usepackage[all]{xy}

% We use 2cell for 2-commutative diagrams.
\xyoption{2cell}
% Theorem environments.
%
\theoremstyle{plain}
\newtheorem{theorem}[subsection]{Theorem}
\newtheorem{proposition}[subsection]{Proposition}
\newtheorem{lemma}[subsection]{Lemma}

\theoremstyle{definition}
\newtheorem{definition}[subsection]{Definition}
\newtheorem{example}[subsection]{Example}
\newtheorem{exercise}[subsection]{Exercise}
\newtheorem{situation}[subsection]{Situation}

\theoremstyle{remark}
\newtheorem{remark}[subsection]{Remark}
\newtheorem{remarks}[subsection]{Remarks}

\numberwithin{equation}{subsection}

% Macros
%
\def\lim{\mathop{\mathrm{lim}}\nolimits}
\def\colim{\mathop{\mathrm{colim}}\nolimits}
\def\Spec{\mathop{\mathrm{Spec}}}
\def\Hom{\mathop{\mathrm{Hom}}\nolimits}
\def\Ext{\mathop{\mathrm{Ext}}\nolimits}
\def\SheafHom{\mathop{\mathcal{H}\!\mathit{om}}\nolimits}
\def\SheafExt{\mathop{\mathcal{E}\!\mathit{xt}}\nolimits}
\def\Sch{\mathit{Sch}}
\def\Mor{\mathop{\mathrm{Mor}}\nolimits}
\def\Ob{\mathop{\mathrm{Ob}}\nolimits}
\def\Sh{\mathop{\mathit{Sh}}\nolimits}
\def\NL{\mathop{N\!L}\nolimits}
\def\CH{\mathop{\mathrm{CH}}\nolimits}
\def\proetale{{pro\text{-}\acute{e}tale}}
\def\etale{{\acute{e}tale}}
\def\QCoh{\mathit{QCoh}}
\def\Ker{\mathop{\mathrm{Ker}}}
\def\Im{\mathop{\mathrm{Im}}}
\def\Coker{\mathop{\mathrm{Coker}}}
\def\Coim{\mathop{\mathrm{Coim}}}

% Boxtimes
%
\DeclareMathSymbol{\boxtimes}{\mathbin}{AMSa}{"02}

%
% Macros for moduli stacks/spaces
%
\def\QCohstack{\mathcal{QC}\!\mathit{oh}}
\def\Cohstack{\mathcal{C}\!\mathit{oh}}
\def\Spacesstack{\mathcal{S}\!\mathit{paces}}
\def\Quotfunctor{\mathrm{Quot}}
\def\Hilbfunctor{\mathrm{Hilb}}
\def\Curvesstack{\mathcal{C}\!\mathit{urves}}
\def\Polarizedstack{\mathcal{P}\!\mathit{olarized}}
\def\Complexesstack{\mathcal{C}\!\mathit{omplexes}}
% \Pic is the operator that assigns to X its picard group, usage \Pic(X)
% \Picardstack_{X/B} denotes the Picard stack of X over B
% \Picardfunctor_{X/B} denotes the Picard functor of X over B
\def\Pic{\mathop{\mathrm{Pic}}\nolimits}
\def\Picardstack{\mathcal{P}\!\mathit{ic}}
\def\Picardfunctor{\mathrm{Pic}}
\def\Deformationcategory{\mathcal{D}\!\mathit{ef}}

% --- appended 2026-10-12 (Chief): the project own definitions of the reference/slogan/history
% display environments, copied mechanically from the shared Stacks preamble lines 8-12
% (\usepackage{verbatim} plus the three \newenvironment lines, which the chapter preamble keeps
% verbatim from the project preamble). Nothing added, removed or redefined: these environments
% carry no mathematical content and the project itself typesets them as comments. Without these
% three lines a counted statement whose own text uses them fails with Environment slogan undefined.
\usepackage{verbatim}
\newenvironment{reference}{\comment}{\endcomment}
\newenvironment{slogan}{\comment}{\endcomment}
\newenvironment{history}{\comment}{\endcomment}
\usepackage{cleveref}
\crefname{originaldefinition}{original definition}{original definitions}
\Crefname{originaldefinition}{Original definition}{Original definitions}
\makeatletter
\newlabel{definition-invertible-sheaf-effective-Cartier-divisor}{{01WX}{source}{}{sourcepointer-1}{}}
\newlabel{definition-invertible-sheaf-effective-Cartier-divisor@cref}{{[originaldefinition][0][]01WX}{[1][0][]source}{}{}{}}
\makeatother
\crefname{originaldefinition}{original definition}{original definitions}
\Crefname{originaldefinition}{Original definition}{Original definitions}
\makeatletter
\newlabel{definition-regular-meromorphic-ideal-denominators}{{02P1}{source}{}{sourcepointer-2}{}}
\newlabel{definition-regular-meromorphic-ideal-denominators@cref}{{[originaldefinition][0][]02P1}{[1][0][]source}{}{}{}}
\makeatother
\crefname{originallemma}{original lemma}{original lemmas}
\Crefname{originallemma}{Original lemma}{Original lemmas}
\makeatletter
\newlabel{lemma-blow-up-pullback-effective-Cartier}{{0809}{source}{}{sourcepointer-3}{}}
\newlabel{lemma-blow-up-pullback-effective-Cartier@cref}{{[originallemma][0][]0809}{[1][0][]source}{}{}{}}
\makeatother
\crefname{originallemma}{original lemma}{original lemmas}
\Crefname{originallemma}{Original lemma}{Original lemmas}
\makeatletter
\newlabel{lemma-blowing-up-affine}{{0804}{source}{}{sourcepointer-4}{}}
\newlabel{lemma-blowing-up-affine@cref}{{[originallemma][0][]0804}{[1][0][]source}{}{}{}}
\makeatother
\crefname{originallemma}{original lemma}{original lemmas}
\Crefname{originallemma}{Original lemma}{Original lemmas}
\makeatletter
\newlabel{lemma-meromorphic-sections-pullback}{{02OY}{source}{}{sourcepointer-5}{}}
\newlabel{lemma-meromorphic-sections-pullback@cref}{{[originallemma][0][]02OY}{[1][0][]source}{}{}{}}
\makeatother
\title{Blowing up the ideal of denominators of a regular meromorphic section makes that section a global regular section modulo the exceptional divisor}
\author{The Stacks Project Authors}
\date{stacks-project:divisors.tex:local-snapshot-20261012; GNU FDL 1.2 or later}
\begin{document}
\maketitle
\noindent\textit{Editorial note.} All delivered mathematical text is the author's own, copied from the pinned Stacks Project chapter file pinned-original-source.tex; no statement or proof step was rewritten, reordered or removed.\par\medskip
\noindent\textit{Editorial note.} Following the multi-claim convention for this source, source\_url is the source-level URL https://stacks.math.columbia.edu/ and the source label lemma-blowing-up-denominators is carried in claim\_key and here rather than in a per-claim tag URL.\par\medskip
\noindent\textit{Editorial note.} The preamble is the chapter's shared Stacks Project preamble, retained as sources/stacks-divisors/retained-preamble.tex: a line-for-line excerpt of the shared Stacks preamble staged with this snapshot. No macro was invented or redefined, and this card adds no preamble text of its own.\par\medskip
\noindent\textit{Editorial note.} Statement span 9232-9248 and proof span 9250-9286 come from the halfA pool card operations/pool-cards/stacks-divisors-halfA.json and were re-read in the pinned chapter file before this card was written.\par\medskip
\noindent\textit{Editorial note.} The proof refers to two definitions of the chapter (the denominator ideal and the invertible sheaf of a divisor difference); both references are part of the counted span and are handled as navigation pointers.\par\medskip
\noindent\textit{Editorial note.} The closing parenthesis giving the divisor cut out by the numerator is the author's and is preserved.\par\medskip
\noindent\textit{Editorial note.} 5 cross-references occur inside the counted spans; they are listed in external\_references as navigation pointers only.\par\medskip
\noindent\textit{Allowed standard prerequisites (editorial).}\begin{itemize}
\item The definition of the ideal of denominators of a regular meromorphic section and its properties, cited from the same chapter. Cited at proof lines 9250-9252. Gives admissibility of the blowing up and lets the section be written locally as a quotient of nonzerodivisors.
\item The notation for the invertible sheaf associated with an effective Cartier divisor and for the twisting of a line bundle by the difference of two divisors, cited from the same chapter. Cited at proof lines 9253-9256. Makes the statement meaningful and converts it into a statement about a regular section of a twisted bundle.
\item Pullbacks of meromorphic sections along a blowing up and along a morphism with integral fibres, cited from the same chapter. Cited at proof lines 9257-9260. Shows that the pullback of the given meromorphic section is defined and nonvanishing in the appropriate sense.
\item The affine local description of a blowing up by blowup algebras and the identification of the exceptional divisor over each chart with the divisor cut out by the chosen element. Cited at proof lines 9272-9276. Supplies the charts and the generator of the twist used in the local computation.
\end{itemize}\medskip
\ifdefined\hypertarget\hypertarget{sourcepointer-1}{}\fi
\noindent\textit{Source reference.} Original definition 01WX: focused-seeds/AGP/sources/stacks-snapshot-20261011/tags.txt, source lines 2498--2498. Same-chapter reference used in the retained author text at source line(s) 3576, 6928, 9257. It is supplied only as a navigation pointer so that the retained text compiles; the counted item delivers the author's own statement and the author's complete proof, and no part of the referenced item is copied in or substituted for the counted argument.\par\medskip
\ifdefined\hypertarget\hypertarget{sourcepointer-2}{}\fi
\noindent\textit{Source reference.} Original definition 02P1: focused-seeds/AGP/sources/stacks-snapshot-20261011/tags.txt, source lines 3536--3536. Same-chapter reference used in the retained author text at source line(s) 9253. It is supplied only as a navigation pointer so that the retained text compiles; the counted item delivers the author's own statement and the author's complete proof, and no part of the referenced item is copied in or substituted for the counted argument.\par\medskip
\ifdefined\hypertarget\hypertarget{sourcepointer-3}{}\fi
\noindent\textit{Source reference.} Original lemma 0809: focused-seeds/AGP/sources/stacks-snapshot-20261011/tags.txt, source lines 10174--10174. Same-chapter reference used in the retained author text at source line(s) 8602, 9259. It is supplied only as a navigation pointer so that the retained text compiles; the counted item delivers the author's own statement and the author's complete proof, and no part of the referenced item is copied in or substituted for the counted argument.\par\medskip
\ifdefined\hypertarget\hypertarget{sourcepointer-4}{}\fi
\noindent\textit{Source reference.} Original lemma 0804: focused-seeds/AGP/sources/stacks-snapshot-20261011/tags.txt, source lines 10169--10169. Same-chapter reference used in the retained author text at source line(s) 8406, 8524, 8536, 8578, 8731, 8837, 8956, 9203, 9275. It is supplied only as a navigation pointer so that the retained text compiles; the counted item delivers the author's own statement and the author's complete proof, and no part of the referenced item is copied in or substituted for the counted argument.\par\medskip
\ifdefined\hypertarget\hypertarget{sourcepointer-5}{}\fi
\noindent\textit{Source reference.} Original lemma 02OY: focused-seeds/AGP/sources/stacks-snapshot-20261011/tags.txt, source lines 3533--3533. Same-chapter reference used in the retained author text at source line(s) 9260. It is supplied only as a navigation pointer so that the retained text compiles; the counted item delivers the author's own statement and the author's complete proof, and no part of the referenced item is copied in or substituted for the counted argument.\par\medskip
\begin{lemma}
\label{lemma-blowing-up-denominators}
Let $X$ be a locally Noetherian scheme.
Let $\mathcal{L}$ be an invertible $\mathcal{O}_X$-module.
Let $s$ be a regular meromorphic section of $\mathcal{L}$.
Let $U \subset X$ be the maximal open subscheme such that
$s$ corresponds to a section of $\mathcal{L}$ over $U$.
The blowup $b : X' \to X$ in the ideal of denominators
of $s$ is $U$-admissible. There exists an effective Cartier divisor
$D \subset X'$ and an isomorphism
$$
b^*\mathcal{L} = \mathcal{O}_{X'}(D - E),
$$
where $E \subset X'$ is the exceptional divisor such that the
meromorphic section $b^*s$ corresponds, via the isomorphism,
to the meromorphic section $1_D \otimes (1_E)^{-1}$.
\end{lemma}

\begin{proof}
From the definition of the ideal of denominators in
Definition
\ref{definition-regular-meromorphic-ideal-denominators}
we immediately see that $b$ is a $U$-admissible blowup.
For the notation $1_D$, $1_E$, and $\mathcal{O}_{X'}(D - E)$
please see Definition
\ref{definition-invertible-sheaf-effective-Cartier-divisor}.
The pullback $b^*s$ is defined by
Lemmas \ref{lemma-blow-up-pullback-effective-Cartier} and
\ref{lemma-meromorphic-sections-pullback}.
Thus the statement of the lemma makes sense.
We can reinterpret the final assertion as saying
that $b^*s$ is a global regular section of
$b^*\mathcal{L}(E)$ whose zero scheme is $D$.
This uniquely defines $D$ hence
to prove the lemma we may work affine locally on $X$ and $X'$.
Assume $X = \Spec(A)$ is affine and
$\mathcal{L} = \mathcal{O}_X$. Then $s$ is a regular meromorphic
function and shrinking further we may assume
$s = a'/a$ with $a', a \in A$ nonzerodivisors.
Then the ideal of denominators of $s$ corresponds
to the ideal $I = \{x \in A \mid xa' \in aA\}$.
Recall that $X'$ is covered by spectra of affine blowup
algebras $A' = A[\frac{I}{x}]$ with $x \in I$
(Lemma \ref{lemma-blowing-up-affine}). Fix $x \in I$ and
write $xa' = a a''$ for some $a'' \in A$.
The divisor $E \subset X'$ is cut out by $x \in A'$ over
the spectrum of $A'$ and hence $1/x$ is a 
generator of $\mathcal{O}_{X'}(E)$ over $\Spec(A')$.
Finally, in the total quotient ring
of $A'$ we have $a'/a = a''/x$. Hence $b^*s = a'/a$ restricts
to a regular section of $\mathcal{O}_{X'}(E)$ which is
over $\Spec(A')$ given by $a''/x$. This finishes the proof.
(The divisor $D \cap \Spec(A')$ is cut out by the image of
$a''$ in $A'$.)
\end{proof}

\end{document}
